\scanpage{64}
Faisons les calculs pour $J=\{0,1,\infty\}\subset\mathbb P^1_S$, on trouve facilement :
\[
\left\{\begin{aligned}
\sigma_0^t(z)&=\frac{z-t}{z-1} &&(\text{échange }0\text{ et }t,\ 1\text{ et }\infty)\\
\sigma_1^t(z)&=t/z &&(\text{échange }1\text{ et }t,\ 0\text{ et }\infty)\\
\sigma_\infty^t(z)&=t\frac{z-1}{z-t} &&(\text{échange }\infty\text{ et }t,\ 0\text{ et }1)
\end{aligned}\right.
\tag*{(34)}
\]
\margin{NB $\sigma_\infty^t=\sigma_0^t\sigma_1^t$}

On trouve donc
\[
\left\{\begin{aligned}
\Delta_0^t &: \text{ solutions de }z^2-2z+t=0,\quad &\Delta_0^t&=\{1\pm\sqrt{1-t}\}\\
\Delta_1^t &: \rule{3em}{.35pt}\quad z^2-t=0 &\Delta_1^t&=\{\pm\sqrt t\}\\
\Delta_\infty^t &: \rule{3em}{.35pt}\quad z^2-2tz+t=0 &\Delta_\infty^t&=\{t\pm\sqrt{t(1-t)}\}
\end{aligned}\right.
\tag*{(35)}
\]
et on doit avoir
\[
\begin{aligned}
f_2^t:\quad &\Delta_0^t\longrightarrow 0\\
&\Delta_1^t\unsure{\text{---}|\longrightarrow}1\\
&\Delta_\infty^t\longrightarrow\infty
\end{aligned}
\qquad\text{\uline{avec multiplicité 2}}
\tag*{(36)}
\]
ce qui nous donne (puisque \deleted{$\unsure{f_2|f}$}les zéros et pôles de $f_2^t$ de degré quatre sont connus)
\[
f_2^t(z)=c_t\left(\frac{z^2-2z+t}{z^2-2tz+t}\right)^2
\qquad c_t\ \text{section de }\underline O_S^*,
\]
dont la valeur s'obtient en écrivant que $\Delta_1^t\longrightarrow1$. Mais si $\unsure{\xi|\zeta}$ est une section de $\Delta_1^t$ i.e. $\xi^2=t$, on a
\[
\begin{aligned}
\xi^2\deleted{+t}-2\xi+t&=2(t-\xi)\\
\xi^2-2t\xi+t&=2t(1-\xi)
\end{aligned}
\]
\[
f_2^t(\xi)=c_t\left(\frac{2(t-\xi)}{2t(1-\xi)}\right)^2
=c_t/t^2\,\deleted{\frac{\xi^2-2t\xi+t^2}{\xi^2-2\xi+1}=}\left(\frac{t-\xi}{1-\xi}\right)^2
\]
Or
\[
\begin{aligned}
\deleted{=\frac{c_t}{t^2}\frac{t^2+t-2t\xi}{1+t-2\xi}}-\quad
\frac{t-\xi}{1-\xi}&=\frac{(t-\xi)(1+\xi)}{(1-\xi)(1+\xi)}
=\frac{-\xi^2+(t-1)\xi+t}{1-\xi^2}\\
&=\frac{(t-1)\xi}{1-t}=-\xi
\end{aligned}
\]
donc
\[
 f_2^t(\xi)=\frac{c_t}{t^2}\xi^2=\frac{c_t t}{t^2}=\frac{c_t}{t}
 \qquad\text{d'où}\qquad c_t=t
\]
% archive sheet number: 63
